% Solution to Problem 5 — O-slice filtration for N_infty operads (Blumberg)
% v8, corrected against First Proof paper (Blumberg-Hill-Lawson, "Generalized equivariant slice categories")

\documentclass[11pt]{article}

\usepackage[margin=1in]{geometry}
\usepackage{amsmath, amssymb, amsthm}
\usepackage{mathtools}
\usepackage{xcolor}
\usepackage{hyperref}
\hypersetup{colorlinks=true,
            linkcolor=blue!50!black,
            citecolor=blue!50!black,
            urlcolor=blue!50!black}

\newcommand{\aO}{\mathcal{O}}
\newcommand{\Sp}{\mathrm{Sp}}
\newcommand{\Z}{\mathbb{Z}}
\newcommand{\R}{\mathbb{R}}
\newcommand{\gconn}{\operatorname{gconn}}
\newcommand{\conn}{\operatorname{conn}}
\newcommand{\Sub}{\operatorname{Sub}}
\DeclareMathOperator{\Map}{Map}

\newtheorem{theorem}{Theorem}
\newtheorem{lemma}{Lemma}
\newtheorem{proposition}{Proposition}
\newtheorem{corollary}{Corollary}
\theoremstyle{definition}
\newtheorem{definition}{Definition}
\newtheorem{remark}{Remark}

\title{Solution to Problem 5: The $\aO$-Slice Filtration for $N_\infty$ Operads}
\author{}
\date{May 27, 2026}

\begin{document}
\maketitle

\section*{Setup: transfer systems and indexing systems}

Fix a finite group $G$. Following Balchin--Barnes--Roitzheim \cite{BBR} and
Rubin \cite{Rubin}, the homotopy theory of $N_\infty$ operads on $G$ is governed
by the combinatorics of \emph{transfer systems}.

\begin{definition}[transfer system]
A \emph{transfer system} on $G$ is a partial order, denoted $\to$, on the set
$\Sub(G)$ of subgroups of $G$ satisfying:
\begin{enumerate}\setlength\itemsep{2pt}
\item it refines subgroup inclusion: if $H\to K$ then $H\subseteq K$;
\item it is conjugation invariant: if $H\to K$ and $g\in G$ then
$gHg^{-1}\to gKg^{-1}$;
\item it is closed under restriction: if $H\to K$ and $J\subseteq K$ then
$H\cap J\to J$.
\end{enumerate}
Transfer systems on $G$ form a poset under refinement.
\end{definition}

\begin{definition}[admissible sets, indexing system]
Let $\aO$ be a transfer system on $G$. A finite $H$-set $T=\coprod_i H/K_i$ is
\emph{admissible} for $\aO$ if $K_i\to H$ for all $i$. The collection $\aO(H)$ of
admissible $H$-sets, as $H$ varies, is the associated \emph{indexing system}; we
also denote the corresponding $N_\infty$ operad by $\aO$.
\end{definition}

For a finite $G$-set $T$ we have the \emph{$T$-norm}
$N^T:\Sp^G\to\Sp^G$, built from the Hill--Hopkins--Ravenel norms \cite{HHR16}
by $N^{G/H}(E)=N^G_H i_H^\ast(E)$ and $N^{T_0\sqcup T_1}(E)=N^{T_0}(E)\otimes
N^{T_1}(E)$.

\section*{The $\aO$-slice filtration}

We use equivariant localization (nullification) to construct the relative slice
tower. Recall that the equivariant localizing subcategory generated by a set of
objects is its closure under homotopy colimits, retracts, and tensors with orbit
spectra.

\begin{definition}[$\aO$-slice connective spectra]
For an indexing system $\aO$, let $\tau^\aO_{\ge n}$ be the equivariant
localizing subcategory of $\Sp^G$ generated by
\[
\bigl\{\, G_+\otimes_H N^T S^1 \;\bigm|\; T\in\aO(H),\ |T|\ge n \,\bigr\}.
\]
This is the category of \emph{$\aO$-slice $n$-connective} spectra.
\end{definition}

\begin{remark}
There is an equivariant homeomorphism $N^T S^1\cong S^{\R\cdot T}$, the
representation sphere of the permutation representation of $T$. So the
$n$-connective spectra are generated by representation spheres of admissible
sets of cardinality $\ge n$. This replaces the regular-representation slice cells
$G_+\wedge_H S^{m\rho_H}$ of the complete (Hill--Hopkins--Ravenel) filtration:
in the incomplete setting one is only permitted the norms indexed by sets that
$\aO$ actually admits.
\end{remark}

\begin{definition}
For an indexing system $\aO$ and $n\ge 0$:
\begin{itemize}\setlength\itemsep{2pt}
\item the $n$th \emph{$\aO$-slice truncation} $P^n_\aO:\Sp^G_{\ge 0}\to
\Sp^G_{\ge 0}$ is the nullification killing $\tau^\aO_{\ge(n+1)}$;
\item the $n$th \emph{$\aO$-slice cover} $P^\aO_n$ is the fiber of
$\mathrm{Id}\Rightarrow P^{n-1}_\aO$;
\item the $n$th \emph{$\aO$-slice} $P^n_{n,\aO}(E)$ is the fiber of
$P^n_\aO(E)\to P^{n-1}_\aO(E)$.
\end{itemize}
A $G$-spectrum $E$ is an \emph{$\aO$-$n$-slice} if $E\in\tau^\aO_{\ge n}$ and
$E\to P^n_\aO E$ is an equivalence.
\end{definition}

That the tower converges (for $E$ connective, $E\xrightarrow{\sim}\varprojlim
P^n_\aO E$) follows because the Postnikov connectivity of $G_+\otimes_H N^T S^1$
is $|T/H|\to\infty$ with $|T|$. The suspension $\Sigma N^T S^1\simeq N^{T\sqcup
\ast}S^1$ shows $\Sigma:\tau^\aO_{\ge k}\to\tau^\aO_{\ge(k+1)}$, whence the
$\infty$-category of $\aO$-$k$-slices is discrete.

\section*{Geometric fixed points and the characteristic function}

Stable equivalences in $\Sp^G$ are detected on geometric fixed points
$\Phi^H$ for all $H\subseteq G$, so the connectivity of geometric fixed points
is the governing invariant.

\begin{definition}[geometric connectivity]
For $E\in\Sp^G$, let $\gconn(E):\Sub(G)\to\Z\cup\{\pm\infty\}$ be
$\gconn(E)(H):=\conn(\Phi^H E)$.
\end{definition}

\begin{definition}[characteristic function]
For a transfer system $\aO$, define $\chi^\aO:\Sub(G)\to\Sub(G)$ by
\[
\chi^\aO(H)\;:=\;\min\{K\mid K\to H\}\;=\;\bigcap_{K\to H}K .
\]
\end{definition}

The index $[G:\chi^\aO(G)]$ is the maximal cardinality of an admissible $G$-set
with a single orbit; equivalently $[G:\chi^\aO(H)]$ measures how large an
admissible $H$-set $\aO$ permits per orbit. This subgroup, \emph{not} $|H|$
itself, is what enters the connectivity bound.

\section*{Main Theorem}

\begin{theorem}\label{thm:main}
Let $\aO$ be a transfer system on $G$ and let $E$ be a connective $G$-spectrum.
Then $E\in\tau^\aO_{\ge n}$ if and only if for all $H\subseteq G$,
\begin{equation}\label{eq:crit}
[H:\chi^\aO(H)]\cdot\gconn(E)(H)\;\ge\;n .
\end{equation}
\end{theorem}

Equivalently, writing
$n=\min_{H\subseteq G}\,[H:\chi^\aO(H)]\cdot\gconn(E)(H)$, one has
$E\in\tau^\aO_{\ge n}$; this $n$ is the \emph{$\aO$-slice connectivity} of $E$.

When $\aO$ is complete, $\chi^\aO(H)=\{e\}$ for all $H$, $[H:\chi^\aO(H)]=|H|$,
and \eqref{eq:crit} recovers the Hill--Yarnall reformulation \cite{HY18} of the
Hill--Hopkins--Ravenel slice connectivity.

\section*{Proof}

\subsection*{Forward direction (necessity)}

\begin{lemma}\label{lem:forward}
If $E\in\tau^\aO_{\ge n}$ then $[H:\chi^\aO(H)]\cdot\gconn(E)(H)\ge n$ for all
$H\subseteq G$.
\end{lemma}

\begin{proof}
By restriction it suffices to treat $H=G$. Geometric fixed points preserve
homotopy colimits and extensions and annihilate induced spectra from proper
subgroups, so it suffices to check the generators $N^T S^1$ for admissible
$G$-sets $T$ with $|T|\ge n$. Decompose $T=\sum_H n_H\, G/H$. Then
$\Phi^G(N^T S^1)=S^{|T/G|}$ with $|T/G|=\sum_H n_H$, while
\[
|T|=\sum_H n_H[G:H]\ge n .
\]
By definition $[G:\chi^\aO(G)]$ is the maximal $[G:H]$ with $G/H\in\aO(\ast)$
(and all others divide it), so
\[
[G:\chi^\aO(G)]\cdot\sum_H n_H \;\ge\; \sum_H n_H[G:H]\;\ge\; n,
\]
i.e.\ $[G:\chi^\aO(G)]\cdot\gconn(N^T S^1)(G)\ge n$, as required. When
$\chi^\aO(G)=\{e\}$ this is exactly \cite[Thm.~2.5]{HY18}.
\end{proof}

\subsection*{Converse via isotropy separation}

Let $\mathcal P$ be the family of proper subgroups of $G$ and consider the
isotropy separation sequence
\[
E\mathcal P_+\otimes E\;\longrightarrow\; E\;\longrightarrow\;
\widetilde{E\mathcal P}\otimes E .
\]
We argue by downward induction on the subgroup lattice. The first piece is built
from inductions $G/H_+\otimes E$ for $H\in\mathcal P$, controlled by the
inductive hypothesis on restrictions; the second piece is \emph{geometric} and
is handled by the following.

\begin{lemma}[geometric slices]\label{lem:geom}
Let $M$ be a geometric Mackey functor. For any $\aO$,
$\Sigma^k HM$ is a $k\cdot[G:\chi^\aO(G)]$-$\aO$-slice.
\end{lemma}

\begin{proof}
Since $M$ is geometric, for any finite $G$-set $T$ the inclusion of fixed points
$S^{|T/G|}\hookrightarrow N^T S^1$ induces an equivalence
$S^{|T/G|}\otimes HM\xrightarrow{\sim} N^T S^1\otimes HM$. For the lower bound,
choose an admissible $T$ with $|T|$ maximal subject to $|T/G|=k$, namely
$T=k\,G/\chi^\aO(G)$ (as $\chi^\aO(G)$ is the minimal subgroup with
$\chi^\aO(G)\to G$); this gives $\Sigma^k HM\in\tau^\aO_{\ge k[G:\chi^\aO(G)]}$.
For the upper bound, if $T$ is admissible with $|T|>k[G:\chi^\aO(G)]$ then
$|T/G|>k$, so
\[
[N^T S^1,\Sigma^k HM]^G\cong[\Phi^G N^T S^1,\Sigma^k HM(G/G)]
\cong[S^{|T/G|},\Sigma^k HM(G/G)]=0 .
\]
Hence $\Sigma^k HM$ is exactly a $k[G:\chi^\aO(G)]$-slice.
\end{proof}

\begin{lemma}\label{lem:induced}
Let $\mathcal F$ be a family and $\tau$ an equivariant localizing subcategory. If
$i_H^\ast E\in i_H^\ast\tau$ for all $H\in\mathcal F$, then $E\mathcal
F_+\otimes E\in\tau$.
\end{lemma}

\begin{proof}
$E\mathcal F_+\otimes E$ lies in the localizing subcategory generated by
$G/H_+\otimes E\cong G_+\otimes_H i_H^\ast E$ for $H\in\mathcal F$; by hypothesis
each such object is in $\tau$. (Same argument as \cite[Lem.~2.4]{HY18}.)
\end{proof}

\begin{proof}[Proof of Theorem~\ref{thm:main}, converse]
Suppose $E$ is connective with $[H:\chi^\aO(H)]\cdot\gconn(E)(H)\ge n$ for all
$H$. In the isotropy separation sequence, Lemma~\ref{lem:geom} (applied to the
homotopy Mackey functors of the geometric part $\widetilde{E\mathcal P}\otimes
E$) shows its $\aO$-slice connectivity is $\ge n$. By downward induction on the
subgroup lattice, the hypothesis on proper subgroups together with
Lemma~\ref{lem:induced} shows $E\mathcal P_+\otimes E$ also has $\aO$-slice
connectivity $\ge n$. Since localizing subcategories are closed under
extensions, $E\in\tau^\aO_{\ge n}$.
\end{proof}

This is precisely the $\aO$-relative form of \cite[Thm.~2.5]{HY18}; the only
new ingredient is that the multiplier $|H|$ is replaced by $[H:\chi^\aO(H)]$,
reflecting that $\aO$ admits norms only up to index $[H:\chi^\aO(H)]$ per orbit.

\section*{Examples and consistency checks}

\paragraph{Complete $\aO$.} Here every $G/H$ is admissible, $\chi^\aO(H)=\{e\}$,
and \eqref{eq:crit} reads $|H|\cdot\gconn(E)(H)\ge n$ for all $H$ --- the
Hill--Hopkins--Ravenel slice connectivity in the Hill--Yarnall formulation
\cite{HY18, HHR16}.

\paragraph{Trivial $\aO$.} No nontrivial transfers: $\chi^\aO(H)=H$ for all $H$,
so $[H:\chi^\aO(H)]=1$ and \eqref{eq:crit} becomes $\gconn(E)(H)\ge n$ for all
$H$. The only admissible sets are trivial orbits, and the filtration is the
geometric Postnikov tower.

\paragraph{$G=\Z/p$.} The two transfer systems are trivial and complete. For the
complete one, $\chi^\aO(\Z/p)=\{e\}$ gives the bound $p\cdot\gconn(E)(\Z/p)\ge n$
at the top together with $\gconn(E)(e)\ge n$; for the trivial one both indices
are $1$ and the conditions are $\gconn(E)(H)\ge n$. These agree with the
$C_p$-slice computations of \cite{HY18}.

\section*{Scope}

The model-categorical and $\infty$-categorical foundations for the
norm functors and equivariant localization are taken from \cite{HHR16, BBR,
Rubin}; the geometric-Mackey-functor input (Lemma~\ref{lem:geom}) adapts
\cite[\S6]{HillPrimer}. The argument is the Hill--Yarnall template
\cite{HY18} carried out with the indexing-system-restricted norms in place of
the full regular-representation cells, and the characteristic function
$\chi^\aO$ tracking the available norms.

\begin{thebibliography}{9}
\bibitem{HHR16} M.~A.~Hill, M.~J.~Hopkins, D.~C.~Ravenel.
  \emph{On the nonexistence of elements of Kervaire invariant one}.
  Ann.\ of Math.\ (2) \textbf{184} (2016), 1--262.
\bibitem{HY18} M.~A.~Hill, C.~Yarnall.
  \emph{A new formulation of the equivariant slice filtration with applications
  to $C_p$-slices}. Proc.\ Amer.\ Math.\ Soc.\ \textbf{146} (2018), 3605--3614.
\bibitem{HillPrimer} M.~A.~Hill.
  \emph{The equivariant slice filtration: a primer}.
  Homology Homotopy Appl.\ \textbf{14} (2012), 143--166.
\bibitem{BBR} S.~Balchin, D.~Barnes, C.~Roitzheim.
  \emph{$N_\infty$-operads and associahedra}.
  Pacific J.\ Math.\ \textbf{315} (2021), 285--304.
\bibitem{Rubin} J.~Rubin.
  \emph{Detecting Steiner and linear isometries operads}.
  Glasg.\ Math.\ J.\ \textbf{63} (2021), 307--342.
\bibitem{Wilson} D.~Wilson.
  \emph{On categories of slices}. Preprint, arXiv:1711.03472, 2017.
\end{thebibliography}

\end{document}